\section{Pengakuan Penggunaan Alat Berbasis Kecerdasan Buatan (AI)}

{
	\renewcommand{\arraystretch}{1.5}
	
	% Kolom 1: c (Center, selebar angka)
	% Kolom 2: l (Left, selebar nama teks AI)
	% Kolom 3: X (Dinamis, mengambil seluruh sisa lebar kertas untuk paragraf)
	\begin{xltabular}{\textwidth}{| c | l | X |}
		
		% --- 1. CAPTION UTAMA ---
		%\caption{Pengakuan Penggunaan AI Generatif}
		%\label{tab:pengakuan-ai} \\
		\hline
		
		% --- HEADER UTAMA ---
		\textbf{No} & \textbf{Nama Alat AI} & \multicolumn{1}{c|}{\textbf{Penggunaan}} \\ \hline
		\endfirsthead
		
		% --- 2. HEADER HALAMAN LANJUTAN ---
		%\caption*{Lanjutan \autoref{tab:pengakuan-ai}} \\
		\hline
		\textbf{No} & \textbf{Nama Alat AI} & \multicolumn{1}{c|}{\textbf{Penggunaan}} \\ \hline
		\endhead
		
		% --- 3. FOOTER SEMENTARA ---
		\hline
		\multicolumn{3}{|r|}{\textit{Berlanjut ke halaman berikutnya...}} \\ \hline
		\endfoot
		
		% --- 4. FOOTER AKHIR ---
		\hline
		\endlastfoot
		
		% --- ISI TABEL ---
		1 & Consensus AI & Digunakan untuk mempercepat pencarian referensi pada tinjauan dan penelitian terdahulu. \\ \hline
		
		2 & NotebookLM & Membantu memberikan rangkuman pada referensi/\textit{paper} yang didapat untuk memudahkan pemahaman isi referensi. \\ \hline
		
		3 & Gemini & Digunakan untuk membantu memperbaiki tata bahasa penulisan serta memberikan saran penyusunan struktur kalimat pada tugas akhir. \\ \hline
		
		4 & Claude & Membantu menyusun dan melakukan format dokumen agar sesuai dengan kaidah ilmiah pada \LaTeX. \\ \hline
		
	\end{xltabular}
}